\ai{\section{Analytical analysis}\label{sec:analysis}
This section was a subsection of \autoref{sec:rocket2d} until this pass. It is promoted to a
section of its own because it is no longer a short appendix to the equations of motion: it now
works through the full configuration space of the pitching moment --- thrust on or off, thrust
axis aligned or misaligned --- and the static-stability result that hangs off one corner of it.
Every label is unchanged, so references from elsewhere in the document still resolve.

Here ${}^r(\cdot)$ denotes a vector's components resolved along the body axes $\hat{x}_r,
\hat{z}_r$ at the current instant --- i.e.\ $m\,{}^r\ddot{x} = (m\ddot{\vec{x}})\cdot\hat{x}_r$
--- \emph{not} a derivative taken within the rotating frame $\mathcal{F}_r$ itself (which would
additionally need that frame's own rotation-coupling terms). The equations below are therefore
just $\sum\vec{F}\cdot\hat{x}_r$, $\sum\vec{F}\cdot\hat{z}_r$ of \autoref{eq:eom_trans},
expanded, and \autoref{eq:eom_rot} with the geometry of \autoref{sec:geometry} substituted in.

Two model-level extensions of \autoref{sec:models} appear in the equations below and are worth
naming up front, because both are \emph{off by default} and both reduce this section to its
previous form when they are:
\begin{itemize}
    \item the thrust cant angle $\delta$ (\autoref{sec:forces}, \autoref{sec:thrust}), which is
    zero for an ideally aligned nozzle;
    \item the gravity-direction rotation angle $\varphi = x/R_e$ (\autoref{sec:gravity} >
    Direction), which is zero at the launch point and stays zero for the whole flight if the
    local vertical is frozen there.
\end{itemize}
Setting $\delta = \varphi = 0$ recovers every equation of this section as it read before those
two models existed.}

\subsection{Forces}\label{sec:analysis_forces}
\ai{The aerodynamic force is a body-axis quantity already; thrust is one too once its own axis is
fixed relative to the body (\autoref{sec:forces}), and only weight has to be resolved. With
$\theta$ measured from the inertial horizontal, positive nose-up (\autoref{sec:frames}), the body
axes read
\begin{equation} \label{eq:body_axes_in_ground}
    \hat{x}_r = \cos\theta\,\hat{x}_{g0} - \sin\theta\,\hat{z}_{g0}, \qquad
    \hat{z}_r = \sin\theta\,\hat{x}_{g0} + \cos\theta\,\hat{z}_{g0},
\end{equation}
where $\hat{x}_{g0}, \hat{z}_{g0}$ are the \emph{launch-site} gravity-frame axes --- the frame
$\theta$ is measured in, and the frame the simulator integrates $x, z$ in
(\autoref{sec:gravity} > Direction). Weight acts along the \emph{local} vertical
$\hat{z}_g(\varphi) = -\sin\varphi\,\hat{x}_{g0} + \cos\varphi\,\hat{z}_{g0}$, so
\begin{equation}
    \vec{W}\cdot\hat{x}_r = -mg\sin(\theta+\varphi), \qquad
    \vec{W}\cdot\hat{z}_r = +mg\cos(\theta+\varphi),
\end{equation}
i.e.\ the downrange rotation of gravity enters the body-axis force balance as nothing more than
a shift of the pitch angle, $\theta \to \theta + \varphi$. That is the whole of its effect, and
it is worth reading once: $\theta+\varphi$ is exactly the rocket's pitch above the \emph{local}
horizon, which dips below the launch-site horizontal by $\varphi$ at downrange distance
$x = R_e\varphi$. Adding thrust and the aerodynamic force,
\begin{align}
    m{}^r\ddot{x} &= -mg\sin(\theta+\varphi) + T\cos\delta - A_x \label{eq:analytic_force_x}\\
    m{}^r\ddot{z} &= mg\cos(\theta+\varphi) + T\sin\delta - A_z \label{eq:analytic_force_z}
\end{align}
The two limiting attitudes fix the signs, taking $\varphi = \delta = 0$ for the check. Nose
straight up ($\theta = 90^\circ$): weight is purely axial and purely retarding,
$m{}^r\ddot{x} = -mg + T - A_x$, with no normal component at all. Nose horizontal ($\theta = 0$):
$\hat{z}_r \equiv \hat{z}_g$, weight is purely normal, $m{}^r\ddot{z} = mg - A_z$, and does not
oppose the thrust in the slightest.

A third check fixes the sign of $\varphi$. Flying level in the launch frame ($\theta = 0$) at
downrange $x > 0$, \autoref{eq:analytic_force_x} gives $m{}^r\ddot{x} = -mg\sin\varphi < 0$:
gravity has acquired a component that \emph{retards} the downrange motion. It should have. The
Earth's centre sits below the launch point, so from a point downrange of it the pull is tilted
backwards, against the flight direction. At $x = 10\,\mathrm{km}$ this is
$g\varphi \approx 1.5\times10^{-2}\,\mathrm{m/s^2}$, and it grows linearly with downrange
distance.

\textbf{Correction.} \autoref{eq:analytic_force_x} and \autoref{eq:analytic_force_z} previously
carried $\cos\theta$ and $\sin\theta$ the other way round. That pair is correct only for a
$\theta$ measured from the \emph{vertical}; this document, and the simulator's rotation matrix,
measure $\theta$ from the horizontal (\autoref{sec:frames}), which swaps the two. The
vertical-launch check above is the quickest way to see it: the old form gave
$-mg\cos 90^\circ = 0$, i.e.\ a rocket standing on its tail feeling no weight along its own axis.}

\subsection{Moment}\label{sec:analysis_moment}
\ai{Substituting $\vec{BC} = \vec{AC}-\vec{AB}$ and $\vec{BD} = \vec{AD}-\vec{AB}$ from
\autoref{sec:geometry} into \autoref{eq:eom_rot}, with $u_D \equiv 0$,
$\hat{x}_r\times\hat{z}_r = -\hat{y}_r$, and the general thrust vector of \autoref{sec:forces}
$\vec{T} = T\cos\delta\,\hat{x}_r + T\sin\delta\,\hat{z}_r$:
\begin{equation} \label{eq:moment_full}
    I\ddot{\theta} = (u_B - u_C)T\cos\delta + (l_C - l_B)T\sin\delta
                     - u_BA_x + (l_B - l_D)A_z - A_m
\end{equation}
Each term reads off the geometry directly. Thrust acts at $C$, whose displacement from the CoG
is $\vec{BC} = (l_B-l_C)\hat{x}_r + (u_B-u_C)\hat{z}_r$: its lateral offset $u_B-u_C$ gives an
arm to the \emph{axial} thrust component $T\cos\delta$, and its axial offset $l_C-l_B$ gives an
arm to the \emph{normal} thrust component $T\sin\delta$. The aerodynamic force acts at $D$, so
the axial arm $l_B - l_D$ carries its normal component $A_z$ and the lateral arm $u_B$ carries
its axial component $A_x$.

\textbf{Correction.} \autoref{eq:moment_full} previously read $(u_B - u_D)T$ for the thrust term
and $(l_B - l_C)A_z$ for the normal-force term --- the point labels $C$ (thrust application
point) and $D$ (aerodynamic reference point) interchanged between the two terms. This is
independent of the sign fix in \autoref{sec:analysis_forces}: moments are taken about $B$, so
weight never enters \autoref{eq:moment_full} and the two errors have nothing to do with each
other.

\paragraph{Weight contributes nothing, whatever direction it points.} The gravity-direction
rotation $\varphi$ of \autoref{sec:analysis_forces} does not appear in
\autoref{eq:moment_full} at all, and cannot: moments are taken about the centre of gravity $B$,
so $\vec{BB}\times\vec{W} = 0$ for \emph{every} orientation of $\vec{W}$. Rotating the local
vertical is therefore a purely translational correction --- exactly like the Coriolis term of
\autoref{sec:noninertial}, and for the same reason (a body force distributed in proportion to
mass has its resultant at the CoG).

\paragraph{$l_C$ is now a live quantity.} It used to drop out entirely. An on-axis thrust
$\vec{T} = T\hat{x}_r$ applied at a point whose displacement from the CoG has axial part
$(l_B-l_C)\hat{x}_r$ contributes $(l_B-l_C)T\,\hat{x}_r\times\hat{x}_r = 0$ from that part:
sliding the application point forward or aft along the axis changed nothing, because the force
was parallel to the displacement it would act through. Only the lateral offset $u_C$ survived.
Once $\delta \neq 0$ the thrust has a component normal to $\hat{x}_r$, and a normal force
\emph{does} care how far along the axis it is applied --- hence the second term of
\autoref{eq:moment_full}, and hence the axial position of the nozzle mattering for the first
time in this document.

\paragraph{The configuration space.} \autoref{eq:moment_full} specialises four ways, and the
remaining subsections take them in turn. Throughout, \emph{aligned} is shorthand for the
on-axis geometry $u_B = u_C = 0$ \emph{and} $\delta = 0$; the simulator has no $u_B$/$u_C$
parameters at all, so $u_B = u_C = 0$ holds unconditionally there and $\delta$ is the only
misalignment it can represent.
\begin{center}
\begin{tabular}{l l l}
 & \textbf{aligned} ($\delta = 0$) & \textbf{misaligned} ($\delta \neq 0$) \\
\textbf{no thrust} ($T=0$) & \autoref{eq:moment_nothrust} & identical --- $T=0$ scales $\delta$ away \\
\textbf{with thrust} ($T>0$) & \autoref{eq:moment_aligned} & \autoref{eq:moment_misaligned} \\
\end{tabular}
\end{center}
Only three of the four cells are distinct, and only one of them is new.}

\subsection{No thrust: coast and descent}\label{sec:analysis_nothrust}
\ai{The coast and descent phases are worth isolating from the boost phase. Setting $T = 0$ drops
the thrust term and leaves the airframe's purely aerodynamic behaviour, free of any dependence
on the motor's thrust level or on where the thrust acts --- so whatever stability it shows is a
property of the geometry and the aerodynamic coefficients alone, not of the motor:
\begin{equation} \label{eq:moment_nothrust}
    I\ddot{\theta} = -u_BA_x + (l_B - l_D)A_z - A_m.
\end{equation}
Thrust misalignment is physically irrelevant in this case rather than merely neglected: $T = 0$
multiplies both $\cos\delta$ and $\sin\delta$ out of \autoref{eq:moment_full}, so the two
right-hand columns of the table above collapse onto one another exactly, for any $\delta$
whatever. It is stated rather than passed over silently because it is the reason the
static-stability result below --- which is built on this case --- needs no revision when the
misalignment model is switched on.}

\subsection{Static stability}\label{sec:analysis_stability}
\paragraph{Static stability} is a steady measure of stability of rocket due to disturbance, which
for this case is angle of attack $\alpha$. We will use no thrust, axis aligned
\ai{--- \autoref{eq:moment_nothrust} with $u_B = 0$, which is the same right-hand side as
\autoref{eq:moment_aligned} below (\autoref{sec:analysis_moment}: with $u_B = u_C = \delta = 0$
the thrust contributes nothing, so the boost and coast phases share one pitching-moment
equation)}. Before the analysis we need to expand and linearize $A_z, A_m$, which are given by
\ai{\begin{equation} \label{eq:aero_linearised}
    A_z = qc_zS,\quad A_m = qc_mSd_{ref},
\end{equation}}
where $q$ is dynamic pressure, $c_z, c_m$ are respectively lift and pitching moment coefficeints and \ai{$S, d_{ref}$} reference area and length, More on that in \ai{\autoref{sec:aero}, where the same two coefficients are written $C_N$ (normal force) and $C_m$ and the reference area and length are the same $S$, $d_{ref}$ (\autoref{eq:aero_magnitudes}).

\textbf{Rename.} The reference area and length were written $A$ and $l$ here until this pass.
They are now $S$ and $d_{ref}$ throughout: it matches \autoref{sec:aero} and the simulator, and
it frees the letter $A$ for the two unrelated things it already denotes in this document --- the
nose-tip point $A$ of \autoref{sec:geometry} and the force magnitudes $A_x, A_z, A_m$. The
dynamic pressure $q$ comes from \autoref{sec:atmosphere} and \autoref{sec:relwind}.}

Coefficents $c_z, c_m$ can be linearized by using
\begin{align}
    c_z &= c_{z,0} + \alpha c_{z,\alpha}, \\
    c_m &= c_{m,0} + \alpha c_{m,\alpha}.
\end{align}

By taking partial derivative of angular acceleration by $\alpha$, we expect $\frac{\partial \ddot\theta}{\partial \alpha} < 0$ behavior, meaning the rocket will stabilize the trajectory. The expansion and derivative gives
\ai{\begin{equation} \label{eq:dtheta_dalpha}
    \frac{\partial \ddot\theta}{\partial \alpha} = \frac{qS}{I}\biggl[(l_B - l_D)c_{z,\alpha} - c_{m,\alpha}d_{ref}\biggr] < 0,
\end{equation}
which is met, provided $c_{z,\alpha} > 0$ (the normal force grows with angle of attack --- the
normal case, and the only one in which dividing through by $c_{z,\alpha}$ preserves the sense of
the inequality), when the stability margin $\textrm{SM}$ satisfies
\begin{equation} \label{eq:static_margin}
    \textrm{SM} = l_B - l_D < d_{ref}\frac{c_{m,\alpha}}{c_{z,\alpha}}.
\end{equation}
Both lengths are measured from the nose tip $A$ (\autoref{sec:geometry}), so $\textrm{SM} = l_B - l_D$ is
\emph{negative} exactly when the aerodynamic reference point $D$ sits aft of the centre of
gravity $B$ --- the familiar ``centre of pressure behind centre of gravity'' rule, with the sign
carried by the ordering of the two lengths rather than by a separate convention. In the limiting
case where $D$ is the centre of pressure itself, $c_{m,\alpha} = 0$ by definition and
\autoref{eq:static_margin} collapses to precisely that rule. The simulator works with the
positive distance $l_1 = l_D - l_B$ instead, i.e.\ $\textrm{SM} = -l_1$
(\autoref{sec:aero}).

\autoref{eq:static_margin} is dimensional. Quoting the margin non-dimensionally --- in
calibers, i.e.\ divided by $d_{ref}$ --- is the rocketry convention and is left for later,
\autoref{sec:backlog}.

\textbf{Correction.} \autoref{eq:dtheta_dalpha} and \autoref{eq:static_margin} previously carried
$l_B - l_C$, inherited from \autoref{eq:moment_aligned}'s interchanged point labels. The static
margin is a centre-of-gravity-to-centre-of-pressure distance and cannot involve the thrust point
$C$ at all; the substitution of $l_D$ for $l_C$ is the only change, and the form of every result
in this subsection is otherwise unaffected.

\textbf{Scope.} This is a \emph{static} criterion: it asks only whether a disturbance in $\alpha$
produces a restoring moment, not what the resulting oscillation does next. The dynamic question
--- natural frequency and damping of the pitch oscillation over a flight whose $q$, $M$, $m$ and
$I$ all change --- needs a pitch-damping derivative this document does not yet have, and is
deferred (\autoref{sec:backlog}).}

\subsection{With thrust, aligned}\label{sec:analysis_aligned}
If thrust and weight axis are aligned then the equation turns into
\ai{\begin{equation} \label{eq:moment_aligned}
    I\ddot{\theta} = (l_B - l_D)A_z - A_m.
\end{equation}
Here ``aligned'' means $u_B = u_C = 0$ \emph{and} $\delta = 0$: with the centre of gravity and
the nozzle both on the body axis and the thrust along it, the first two terms of
\autoref{eq:moment_full} vanish together, the remaining lateral arm $u_B$ in
\autoref{eq:moment_nothrust} vanishes too, and the pitching behaviour reduces to the normal force
acting through the arm $l_B - l_D$ against the airframe's own pitching moment.

The notable feature of this case is what is \emph{missing}: the motor. Thrust has dropped out of
the pitching equation entirely, so an aligned rocket's attitude dynamics are the same during
boost as during coast, however large $T$ is. This is why the simulator computed a thrust moment
of identically zero for every run before the misalignment model existed --- not as an
approximation, but as an exact consequence of the geometry it was given.}

\subsection{With thrust, misaligned}\label{sec:analysis_misaligned}
\ai{With the nozzle canted by a small fixed $\delta$ and the geometry otherwise on-axis
($u_B = u_C = 0$):
\begin{equation} \label{eq:moment_misaligned}
    I\ddot{\theta} = \underbrace{(l_C - l_B)T\sin\delta}_{M_\delta} + (l_B - l_D)A_z - A_m.
\end{equation}
The new term $M_\delta$ is the point of modelling $\delta$ at all, and it behaves quite unlike
everything else in this section. Every other contribution to $\ddot{\theta}$ is aerodynamic and
therefore vanishes with $q$ and with $\alpha$; $M_\delta$ does not. It survives $A_z = A_m = 0$,
it survives $\alpha = 0$, and it survives $q = 0$. For as long as the motor burns, the airframe
carries a built-in disturbance torque proportional to the thrust it is producing --- largest, in
absolute terms, exactly at the thrust peak, and largest \emph{relative} to the restoring moment
just off the pad and again high up, where $q$ is small.

Its sign is worth pinning down once. $l_C > l_B$ for any real rocket (the nozzle is aft of the
centre of gravity), and $\hat{z}_r$ points towards the rocket's belly (\autoref{sec:frames}), so
a positive $\delta$ pushes the tail towards $+\hat{z}_r$ and pitches the nose \emph{up},
$M_\delta > 0$. Both factors are constants of the airframe rather than of the flight state, up
to $T(t)$ and the CoG travel $l_B(t)$ of \autoref{sec:mass}.

\paragraph{It does not change the static-stability result.} $\delta$ is a constant, so
$\partial M_\delta/\partial\alpha = 0$ and \autoref{eq:dtheta_dalpha} --- the derivative that
\autoref{eq:static_margin} is built on --- is untouched. A misalignment adds a constant
\emph{offset} to $\ddot{\theta}$, not a slope: it shifts the equilibrium the airframe settles
at, without changing whether that equilibrium is stable. Concretely, setting $\ddot{\theta} = 0$
in \autoref{eq:moment_misaligned} with the linearised coefficients of
\autoref{eq:aero_linearised} gives the trim angle of attack
\begin{equation} \label{eq:alpha_trim}
    \alpha_{trim} = -\frac{M_\delta}{I\,\partial\ddot{\theta}/\partial\alpha}
        = \frac{-(l_C - l_B)T\sin\delta}
               {qS\bigl[(l_B - l_D)c_{z,\alpha} - c_{m,\alpha}d_{ref}\bigr]},
\end{equation}
taking $c_{z,0} = c_{m,0} = 0$ for a body of revolution. A statically stable airframe has a
negative denominator (\autoref{eq:dtheta_dalpha}), so $\alpha_{trim}$ carries the sign of
$M_\delta$: the rocket does not diverge, it flies slightly nose-off the relative wind, at a
constant small angle whose normal force then curves the trajectory. That curvature, not a
tumble, is the observable consequence of a manufacturing-grade misalignment on a stable rocket.

Two limits are worth reading off \autoref{eq:alpha_trim} rather than discovering in a run.
$\alpha_{trim} \propto T/q$, so it blows up wherever there is thrust without dynamic pressure:
immediately after ignition, and during any high-altitude burn. And it blows up outright as the
static margin approaches zero, since the denominator is $I$ times
\autoref{eq:dtheta_dalpha} --- a marginally stable rocket has no authority to trim against even
a tiny built-in torque. Neither limit is a defect of the model; both are the physical reason
launch rails exist.

\paragraph{What is \emph{not} modelled here.} $\delta$ is a fixed manufacturing-type cant, not a
commanded input: there is no control law, no gimbal rate, no actuator dynamics. And it is the
only misalignment carried --- a lateral offset of the nozzle or of the centre of gravity
($u_B, u_C \neq 0$), which would add the $(u_B-u_C)T\cos\delta$ term of
\autoref{eq:moment_full}, is derived above but has no parameters behind it in the simulator.
Both remain future work, \autoref{sec:backlog}.}
